\documentclass[11pt,a4paper]{article}
\usepackage{local-overleaf}

\begin{document}
\RaggedRight

\makeanimacover

% 2
\kicker{01 · Lectura ejecutiva}
\section{La oportunidad}
\sectionintro{Anima puede transformar una demostración premiada en una capacidad pública concreta: escuchar señales de bienestar, comprender patrones territoriales y orientar respuestas oportunas.}

\begin{navycard}
  \whitecardtitle{Una propuesta nacida en Salta}
  El proyecto obtuvo el primer premio del Hackathon NOA Innova 2026. La comunicación oficial del Municipio destacó un ecosistema digital orientado a la prevención y al abordaje de situaciones vinculadas con la salud mental y el suicidio, con autorización de las personas usuarias y criterios de seguridad y privacidad.
\end{navycard}
\sourceinline{Fuente: Municipalidad de Salta, 10 de agosto de 2026.}{https://prensa.municipalidadsalta.gob.ar/un-proyecto-para-abordar-la-salud-mental-fue-premiado-en-el-hackathon-noa-innova-2026/}

\vspace{5mm}
\begin{minipage}[t]{.31\textwidth}
  \begin{softcard}[height=42mm]
    \cardtitle{Ciudadanía}
    Un registro breve y voluntario que ayuda a reconocer cómo se siente una persona y a acercarle recursos de cuidado.
  \end{softcard}
\end{minipage}\hfill
\begin{minipage}[t]{.31\textwidth}
  \begin{softcard}[height=42mm]
    \cardtitle{Territorio}
    Información agregada por zona y grupo para detectar tendencias sin exponer historias individuales.
  \end{softcard}
\end{minipage}\hfill
\begin{minipage}[t]{.31\textwidth}
  \begin{softcard}[height=42mm]
    \cardtitle{Gestión pública}
    Un observatorio que traduzca señales en prioridades, seguimiento y mejores decisiones de política pública.
  \end{softcard}
\end{minipage}

\vspace{4mm}
\begin{terracard}
  \cardtitle{Punto de partida honesto}
  Hoy existe una demostración visual e interactiva. Sirve para comprender la experiencia y el potencial del producto, pero todavía no constituye un sistema municipal productivo. El siguiente paso responsable no es fijar un precio aislado: es definir, junto con el Municipio, el alcance del primer piloto.
\end{terracard}

\vspace{4mm}
\begin{softcard}
  \cardtitle{Hipótesis a validar}
  Un piloto debe demostrar que una experiencia voluntaria, protegida y medible puede producir información territorial útil sin exponer casos individuales ni sustituir la atención profesional.
\end{softcard}

\clearpage

% 3
\kicker{02 · La aplicación}
\section{Qué es Anima}
\sectionintro{Una experiencia de bienestar emocional que conecta a la persona, un canal de acompañamiento y la lectura agregada para la gestión pública.}

\begin{minipage}[t]{.37\textwidth}
  \vspace{0pt}\centering
  \phoneframe[.54\linewidth]{\AnimaMoodImage}
  \par\vspace{2mm}{\footnotesize\color{Muted}Vista de la demostración actual.}
\end{minipage}\hfill
\begin{minipage}[t]{.59\textwidth}
  \vspace{0pt}
  \begin{card}[fontupper=\small,top=3mm,bottom=3mm]
    \cardtitle{1. Registro y consentimiento}
    La persona expresa cómo se siente mediante una interacción breve y no diagnóstica. Todo contexto adicional requiere una finalidad explícita y consentimiento verificable.
  \end{card}
  \begin{card}[fontupper=\small,top=3mm,bottom=3mm]
    \cardtitle{2. Acompañamiento}
    WhatsApp funciona como puerta de entrada, recordatorio y canal de orientación. No reemplaza a profesionales ni servicios de emergencia.
  \end{card}
  \begin{card}[fontupper=\small,top=3mm,bottom=3mm]
    \cardtitle{3. Observatorio municipal}
    El Municipio visualiza indicadores agregados por zona, período y grupo. El objetivo es decidir dónde reforzar prevención, comunicación o capacidad territorial.
  \end{card}
\end{minipage}

\vspace{4mm}
\begin{softcard}
  \cardtitle{Experiencia ciudadana}
  \begin{minipage}[t]{.31\linewidth}\vspace{0pt}
    \textbf{Interacción breve}\par
    Una experiencia simple y comprensible.
  \end{minipage}\hfill
  \begin{minipage}[t]{.31\linewidth}\vspace{0pt}
    \textbf{Lenguaje de cuidado}\par
    Sin diagnóstico ni etiquetas personales.
  \end{minipage}\hfill
  \begin{minipage}[t]{.31\linewidth}\vspace{0pt}
    \textbf{Participación voluntaria}\par
    La persona decide cuándo registrar cómo se siente.
  \end{minipage}
\end{softcard}

\vspace{4mm}
\begin{navycard}
  \whitecardtitle{Principio rector}
  La plataforma no busca etiquetar ni diagnosticar personas. Busca convertir registros voluntarios y protegidos en conocimiento colectivo útil para el cuidado y la planificación.
\end{navycard}

\clearpage

% 4
\kicker{03 · Valor público}
\section{Cómo se traduce en capacidad municipal}
\sectionintro{La tecnología sólo genera valor cuando conecta una señal con una decisión, una responsabilidad y una forma de evaluar el resultado.}

\begin{center}
\begin{tikzpicture}[node distance=5mm]
  \node[rounded corners=2mm,fill=Sky,minimum width=37mm,minimum height=24mm,align=center,text=Navy,font=\small\bfseries] (a) {Registro\\voluntario};
  \node[rounded corners=2mm,fill=Sand,minimum width=37mm,minimum height=24mm,align=center,text=Navy,font=\small\bfseries,right=of a] (b) {Protección y\\agregación};
  \node[rounded corners=2mm,fill=Sage!20,minimum width=37mm,minimum height=24mm,align=center,text=Navy,font=\small\bfseries,right=of b] (c) {Lectura\\territorial};
  \node[rounded corners=2mm,fill=Terracotta!14,minimum width=37mm,minimum height=24mm,align=center,text=Navy,font=\small\bfseries,right=of c] (d) {Decisión y\\seguimiento};
  \draw[-{Latex[length=2.2mm]},thick,color=Muted] (a)--(b);
  \draw[-{Latex[length=2.2mm]},thick,color=Muted] (b)--(c);
  \draw[-{Latex[length=2.2mm]},thick,color=Muted] (c)--(d);
\end{tikzpicture}
\end{center}

\vspace{3mm}
\begin{minipage}[t]{.48\textwidth}
  \begin{softcard}[height=59mm]
    \cardtitle{Preguntas que el tablero debería responder}
    \begin{itemize}
      \item ¿Qué zonas muestran un cambio sostenido?
      \item ¿Dónde hay baja cobertura de registros?
      \item ¿Qué grupos requieren acciones preventivas?
      \item ¿Qué intervención municipal conviene priorizar?
    \end{itemize}
  \end{softcard}
\end{minipage}\hfill
\begin{minipage}[t]{.48\textwidth}
  \begin{softcard}[height=59mm]
    \cardtitle{Decisiones que podría habilitar}
    \begin{itemize}
      \item Reforzar campañas o talleres en una zona.
      \item Coordinar recursos con CAPS e instituciones.
      \item Evaluar alcance, participación y continuidad.
      \item Detectar vacíos de información antes de actuar.
    \end{itemize}
  \end{softcard}
\end{minipage}

\vspace{5mm}
\begin{terracard}
  \cardtitle{La unidad de valor no es “tener un dashboard”}
  Es poder tomar una decisión concreta con información suficiente, protegida, explicable y actualizada a la frecuencia que la gestión realmente necesita.
\end{terracard}

\vspace{4mm}
\begin{softcard}
  \cardtitle{Cómo se evaluaría el piloto}
  Participación y cobertura territorial · calidad y oportunidad de los datos · decisiones municipales documentadas · cumplimiento de privacidad y respuesta humana.
\end{softcard}

\clearpage

% 5
\kicker{04 · Punto de partida}
\section{De demostración a piloto}
\sectionintro{La demo permite conversar sobre el producto. El piloto debe probar una capacidad operativa acotada, medible y segura.}

\begin{card}
  \begin{minipage}[t]{.21\linewidth}\vspace{0pt}
    {\small\bfseries\color{Sage}ESTADO ACTUAL}\par\smallskip
    {\Large\bfseries\color{Navy}Demostración}
  \end{minipage}\hfill
  \begin{minipage}[t]{.74\linewidth}\vspace{0pt}
    \begin{itemize}
      \item Flujo visual de registro emocional y contacto por WhatsApp.
      \item Dashboard y mapa demostrativos.
      \item Narrativa de consentimiento e integraciones.
      \item Hoja de ruta de producto en cinco etapas.
    \end{itemize}
  \end{minipage}
\end{card}

\vspace{3mm}
\begin{card}
  \begin{minipage}[t]{.21\linewidth}\vspace{0pt}
    {\small\bfseries\color{Terracotta}PRÓXIMO NIVEL}\par\smallskip
    {\Large\bfseries\color{Navy}Piloto}
  \end{minipage}\hfill
  \begin{minipage}[t]{.74\linewidth}\vspace{0pt}
    \begin{itemize}
      \item Definir población, territorio, duración y fecha de inicio.
      \item Construir backend, base de datos e integraciones reales.
      \item Acordar consentimiento, roles, auditoría y retención.
      \item Operar un protocolo humano con capacidad de respuesta.
      \item Medir indicadores y criterios de aceptación.
    \end{itemize}
  \end{minipage}
\end{card}

\vspace{5mm}
\begin{navycard}
  \whitecardtitle{Propuesta para el primer contrato}
  Un alcance deliberadamente acotado: experiencia ciudadana, canal de WhatsApp y observatorio agregado, en una población y un territorio definidos, con acompañamiento clínico y gobernanza de datos. Las integraciones avanzadas, los beneficios y la IA quedan fuera hasta validar los fundamentos.
\end{navycard}

\vspace{5mm}
\begin{center}
\begin{tikzpicture}[node distance=2.2mm]
  \node[rounded corners=2.5mm,fill=Sky,text=Navy,font=\bfseries\scriptsize,minimum width=25mm,minimum height=7mm] (d) {Demo actual};
  \node[rounded corners=2.5mm,fill=Sky,text=Navy,font=\bfseries\scriptsize,minimum width=34mm,minimum height=7mm,right=of d] (a) {Definición conjunta};
  \node[rounded corners=2.5mm,fill=Terracotta!18,text=Navy,font=\bfseries\scriptsize,minimum width=32mm,minimum height=7mm,right=of a] (p) {Piloto controlado};
  \node[rounded corners=2.5mm,fill=Sky,text=Navy,font=\bfseries\scriptsize,minimum width=26mm,minimum height=7mm,right=of p] (e) {Evaluación};
  \node[rounded corners=2.5mm,fill=Sky,text=Navy,font=\bfseries\scriptsize,minimum width=27mm,minimum height=7mm,right=of e] (s) {Decisión de escala};
\end{tikzpicture}
\end{center}

\clearpage

% 6
\kicker{05 · Evolución del producto}
\section{Hoja de ruta propuesta}
\sectionintro{Las cinco etapas originales se conservan, pero se ordenan por riesgo y valor. La primera es la candidata al piloto; las demás se habilitan sólo con evidencia.}

\phase{1}{Núcleo del piloto}{Bot de WhatsApp para convocatoria y acompañamiento; aplicación de registro emocional; observatorio municipal con información agregada y mapa.}
\phase{2}{Historial e integraciones autorizadas}{Calendario personal e integraciones con fuentes de salud o actividad, únicamente con consentimiento, finalidad definida y validación de calidad.}
\phase{3}{Perfiles poblacionales}{Segmentaciones agregadas que ayuden a comprender patrones colectivos, sin convertirlas en etiquetas individuales ni decisiones automáticas.}
\phase{4}{Beneficios saludables}{Incentivos vinculados con acciones de cuidado o acceso a recursos. No se premian crisis, cantidad de registros ni rachas emocionales.}
\phase{5}{IA acotada y supervisada}{Funciones específicas, trazables y evaluadas. Sin diagnóstico, tratamiento ni conversación clínica autónoma.}

\begin{terracard}
  \cardtitle{Puerta de avance}
  Cada etapa debe tener un propósito municipal, datos habilitantes, responsables, criterios de aceptación y una revisión de privacidad y seguridad antes de pasar a la siguiente.
\end{terracard}

\clearpage

% 7
\kicker{06 · Definiciones prioritarias}
\section{Primera ronda de decisiones}
\sectionintro{Estas definiciones permiten transformar la propuesta en un alcance ejecutable, con responsables y criterios verificables.}

\subsection{Alcance y modalidad de trabajo}
\decision{1}{¿La necesidad inmediata es una demostración, un piloto controlado o un sistema municipal productivo?}
\decision{2}{¿Qué etapas deben formar parte del primer contrato?}
\decision{3}{¿Cuál es la ventana esperada de lanzamiento?}
\decision{4}{¿Quién desarrollará y operará cada componente: Equipo Anima, Municipio o un esquema conjunto?}

\subsection{Datos disponibles}
\decision{5}{¿Qué conjuntos de datos existen actualmente y quién es responsable de cada uno?}
\decision{6}{¿Qué nivel de detalle tienen: ciudad, barrio, radio censal, CAPS o domicilio?}
\decision{7}{¿Con qué frecuencia se actualizan?}
\decision{8}{¿En qué formato están: CSV, Excel, base SQL, API o sistema cerrado?}

\clearpage

% 8
\kicker{06 · Definiciones prioritarias}
\section{Segunda ronda de decisiones}
\sectionintro{El observatorio, WhatsApp y los incentivos sólo deben construirse cuando su uso, acceso y propósito estén definidos.}

\subsection{Observatorio y mapa}
\decision{9}{¿Cuál será el mínimo de personas necesario para mostrar una zona sin riesgo de reidentificación?}
\decision{10}{¿Quiénes podrán acceder, filtrar y exportar información?}
\decision{11}{¿Se necesitan distintos roles y permisos?}
\decision{12}{¿El mapa debe actualizarse en tiempo real, diariamente o mensualmente?}
\decision{13}{¿Qué decisión concreta debería poder tomar el Municipio mirando el dashboard?}

\subsection{WhatsApp e incentivos}
\decision{14}{¿La Municipalidad ya cuenta con Meta Business Manager, número verificado y cuenta de WhatsApp Business?}
\decision{15}{¿De dónde surgirá la lista de destinatarios y existe consentimiento explícito para recibir mensajes?}
\decision{16}{¿Qué comportamiento saludable se quiere incentivar?}

\clearpage

% 9
\kicker{07 · Condiciones de confianza}
\section{Privacidad, seguridad y respuesta humana}
\sectionintro{En salud mental, la confianza no es un complemento técnico: es una condición de existencia del proyecto.}

\begin{minipage}[t]{.48\textwidth}
  \begin{card}
    \cardtitle{Gobernanza de datos}
    \begin{itemize}
      \item Minimización: recopilar sólo lo indispensable.
      \item Consentimiento claro, específico y revocable.
      \item Pseudonimización y separación de identidades.
      \item Resultados agregados con umbral mínimo por zona.
      \item Roles, trazabilidad de accesos y exportaciones.
      \item Retención, eliminación e incidentes documentados.
      \item QR de beneficios separado de registros emocionales.
    \end{itemize}
  \end{card}
\end{minipage}\hfill
\begin{minipage}[t]{.48\textwidth}
  \begin{card}
    \cardtitle{Respuesta humana}
    \begin{itemize}
      \item El equipo de psicología define y aprueba el umbral.
      \item Se acuerdan responsable, horario y tiempo de respuesta.
      \item Cada señal tiene una acción posible y un cierre registrable.
      \item Se define qué ocurre fuera de horario o sin capacidad.
      \item Riesgo agudo y emergencias siguen protocolos específicos.
      \item La plataforma no promete una atención que el sistema no pueda brindar.
    \end{itemize}
  \end{card}
\end{minipage}

\vspace{5mm}
\begin{navycard}
  \whitecardtitle{Límite explícito}
  Ningún umbral individual, alerta automática, diagnóstico o recomendación clínica se habilita sólo por decisión de producto. Requiere conducción profesional, protocolo aprobado, capacidad real de respuesta y validación legal.
\end{navycard}

\vspace{4mm}
\begin{terracard}
  \cardtitle{Lenguaje preciso}
  Hasta demostrar que la reidentificación es materialmente imposible, los datos deben describirse como \textbf{pseudonimizados}, no como anónimos.
\end{terracard}

\clearpage

% 10
\kicker{08 · Propuesta económica}
\section{Hacia una cotización defendible}
\sectionintro{Este documento no fija un monto. La propuesta económica se emitirá cuando el Municipio y el Equipo Anima hayan cerrado las definiciones que cambian materialmente el alcance.}

\begin{navycard}
  \whitecardtitle{No se solicita aprobar un presupuesto en esta instancia}
  Se solicita acordar el problema a resolver, la primera etapa, la población, los datos disponibles, el modelo operativo y los criterios de éxito. Con esa base se entregará una cotización formal, comparable y trazable.
\end{navycard}

\vspace{4mm}
\begin{minipage}[t]{.48\textwidth}
  \begin{softcard}
    \cardtitle{La cotización deberá separar}
    \begin{enumerate}
      \item Implementación inicial del piloto.
      \item Requisitos y aportes a cargo del Municipio.
      \item Operación y soporte recurrentes.
      \item Expansiones opcionales posteriores.
    \end{enumerate}
  \end{softcard}
\end{minipage}\hfill
\begin{minipage}[t]{.48\textwidth}
  \begin{softcard}
    \cardtitle{La cotización deberá explicitar}
    \begin{enumerate}
      \item Entregables y exclusiones.
      \item Hitos de aceptación y pagos.
      \item Responsables y dependencias.
      \item Costos recurrentes y vigencia.
    \end{enumerate}
  \end{softcard}
\end{minipage}

\vspace{4mm}
\begin{terracard}
  \cardtitle{Variables que determinan la inversión}
  Alcance funcional · cantidad de personas y zonas · canales e integraciones · estado y calidad de los datos · requisitos de infraestructura · cobertura humana y clínica · soporte esperado · propiedad y entrega de activos · criterios de aceptación.
\end{terracard}

\vspace{8mm}
\begin{softcard}
  \cardtitle{Resultado esperado}
  Una propuesta económica formal para el piloto, con precio, cronograma, supuestos, responsabilidades y opción de continuidad; no una suma preliminar de costos aislados.
\end{softcard}

\vspace{5mm}
\subsection{Secuencia para llegar al monto}
\begin{minipage}[t]{.31\textwidth}
  \begin{card}[height=31mm]
    \numberbadge{1}\hspace{2mm}{\bfseries\color{Navy}Cerrar alcance}\par\smallskip
    {\small Funciones, población, territorio y responsabilidades.}
  \end{card}
\end{minipage}\hfill
\begin{minipage}[t]{.31\textwidth}
  \begin{card}[height=31mm]
    \numberbadge{2}\hspace{2mm}{\bfseries\color{Navy}Cotizar}\par\smallskip
    {\small Entregables, cronograma, supuestos y costos recurrentes.}
  \end{card}
\end{minipage}\hfill
\begin{minipage}[t]{.31\textwidth}
  \begin{card}[height=31mm]
    \numberbadge{3}\hspace{2mm}{\bfseries\color{Navy}Aprobar}\par\smallskip
    {\small Hitos, aceptación, pagos y condiciones de continuidad.}
  \end{card}
\end{minipage}

\clearpage

% 11
\kicker{09 · Cierre}
\section{Próximos pasos}
\sectionintro{El objetivo inmediato es acordar una mesa de alcance y convertir la propuesta en un piloto acotado, seguro y evaluable.}

\phase{1}{Designar contrapartes}{Nombrar un referente institucional, uno técnico y uno clínico para conducir las definiciones.}
\phase{2}{Realizar la mesa de alcance}{Resolver las 16 preguntas prioritarias y documentar decisiones, dudas y responsables.}
\phase{3}{Relevar datos y capacidades}{Confirmar fuentes, formatos, frecuencia, autorizaciones, sistemas existentes y capacidad de respuesta.}
\phase{4}{Diseñar el piloto}{Acordar población, zona, canales, protocolo, indicadores de éxito y criterios de aceptación.}
\phase{5}{Emitir la propuesta formal}{Presentar alcance, cronograma, responsabilidades, inversión y operación recurrente para aprobación.}

\begin{navycard}
  \whitecardtitle{Una decisión concreta para esta reunión}
  Acordar si la Municipalidad desea avanzar con una etapa de definición conjunta orientada a diseñar el primer piloto de Anima en la ciudad de Salta.
\end{navycard}

\vspace{5mm}
\subsection{Fuentes institucionales}
\sourceitem{1. Proyecto ganador del Hackathon NOA Innova 2026}{https://prensa.municipalidadsalta.gob.ar/un-proyecto-para-abordar-la-salud-mental-fue-premiado-en-el-hackathon-noa-innova-2026/}
\sourceitem{2. Apertura del encuentro NOA Innova 2026}{https://prensa.municipalidadsalta.gob.ar/con-mas-de-600-participantes-comenzo-el-noa-innova-hackathon-2026/}

\vspace{10mm}
{\centering\animaWordmark\par
\small\color{Muted}Tecnología cívica para comprender, cuidar y decidir mejor.\par}

\end{document}
